\documentclass[10pt,letterpaper]{article}
\usepackage{cornell-notes}

\title{Chapter 2: Random Subset}
\author{}
\date{}

\setDocTitle{Chapter 2: Random Subset}
\setDocSubtitle{Combinatorial Algorithms}
\setDocVersion{v1.0}
\setDocStatus{Study Companion}
\setDocOwner{Combinatorial Algorithms Cornell Notes}
\setDocAuthor{}
\setDocDate{\today}
\setDocSummary{Cornell-format study and review notes for Chapter 2: Random Subset.}

\setCornellCollection{Combinatorial Algorithms Cornell Notes}
\setCornellUnitType{Chapter}
\setCornellUnitNumber{2}
\setCornellUnitTitle{Random Subset}
\setCornellCompanionLabel{Study and review companion}

\hypersetup{pdftitle={Chapter 2: Random Subset},pdfsubject={Cornell notes},pdfauthor={}}

\begin{document}
\maketitle
\begin{tcolorbox}[colback=Paper,colframe=Ink,boxrule=1pt,arc=2mm]
{\small\bfseries\color{Teal}CORNELL NOTES}\par{\LARGE\bfseries Chapter 2: Random Subset of an $n$-Set (RANSUB)}\par\medskip
\textbf{Topic:} Uniform subset sampling\hfill\textbf{Date:} \today\par
\textbf{Study purpose:} Connect independent membership trials to a uniform distribution on the power set.
\end{tcolorbox}
\begin{tcolorbox}[colback=Teal!5,colframe=Teal,title=Essential question]
Why does choosing every membership bit independently with probability $1/2$ select each of the $2^n$ subsets with the same probability?
\end{tcolorbox}
\begin{tcolorbox}[colback=white,colframe=Mist,title=Learning objectives]
\begin{itemize}\item Implement the RANSUB membership rule.\item Derive the probability of any fixed output subset.\item Identify the required random-number assumption.\item Interpret the sample frequency check without demanding equal observed counts.\end{itemize}
\textbf{Prerequisites:} subsets, independent Bernoulli trials, expectation, and elementary goodness-of-fit reasoning.
\end{tcolorbox}
\section*{Cornell notes}
\begin{longtable}{@{}Q!{\color{Mist}\vrule width 1pt}A@{}}\cornellhead
\cue{What is the output representation?}{An integer array $a_1,\ldots,a_n$ represents a subset: $a_i=1$ when element $i$ is included and $a_i=0$ otherwise. There is no sequencing state because each call produces a fresh independent sample.}
\cue{What is the algorithm?}{For each $i=1,\ldots,n$, draw a random real $\xi_i$ uniformly from $[0,1)$ and set $a_i=\lfloor 2\xi_i\rfloor$. Equivalently, flip a fair coin once per element and include the element on one of the two outcomes.}
\cue{Why is selection uniform?}{Fix any subset $S$. Its bit vector prescribes one outcome for each of $n$ independent fair trials. Therefore
\[
\Pr\{\text{output}=S\}=\prod_{i=1}^{n}\frac12=2^{-n}.
\]
There are $2^n$ possible bit vectors, so the probabilities sum to one.}
\cue{What assumptions matter?}{The random draws must be independent and each threshold event must have probability $1/2$. A biased or correlated generator changes the output distribution. Mapping a discrete generator to bits also must avoid a threshold or modulo bias.}
\cue{Cost and storage}{The method performs one random draw and one assignment per universe element: $n$ trials and an output array of length $n$. \textbf{Derived insight:} the running time is $\Theta(n)$ because the full bit-vector output itself has length $n$.}
\cue{How should sample output be read?}{Repeated uniform sampling does not make every subset appear equally often in a finite run. The supplied experiment samples all subsets of a five-element set many times and evaluates the deviations from the common expected count with a chi-square statistic. Variation around the expected frequency is normal; exact equality would be unusually strict.}
\cue{Historical or legacy context}{The routine obtains reals from an external random-number function and converts them directly to integer bits. In a modern implementation, specify the generator, seed or reproducibility policy, and exact distribution API rather than relying on an implicit global generator.}
\cue{Implementation checklist}{\begin{itemize}\item Produce exactly $n$ bits.\item Use one unbiased independent decision per element.\item Document whether calls are reproducible.\item Test $n=0$ and $n=1$ according to the host API contract.\item Use many trials, not one sample, for distribution checks.\end{itemize}}
\cue{Retrieval practice}{\begin{enumerate}\item What does $a_i$ mean?\item Why is the probability $2^{-n}$?\item Does uniformity imply equal empirical counts?\item Which failure modes introduce bias?\item What is the space used for the output?\end{enumerate}}
\end{longtable}
\begin{tcolorbox}[colback=Ink!4,colframe=Ink,title=Cornell summary]
RANSUB samples a subset by treating every element as an independent yes-or-no decision. The output is the indicator vector of those decisions. Because a particular subset fixes all $n$ bit values, and every required bit occurs with probability $1/2$, its probability is $2^{-n}$. This is the same for every subset, which proves uniformity over the power set. The method is simple and output-linear, but its guarantee depends entirely on unbiased, independent random decisions. A finite experiment should show fluctuations around the expected frequency rather than identical counts, so distribution testing must assess aggregate deviations over many samples. The enduring algorithmic principle is to obtain a uniform choice from a Cartesian product by independently sampling each equally sized coordinate choice.
\end{tcolorbox}
\end{document}
